% GENERATED FILE. Edit canonical lessons, then run npm run generate:books.
% canonical-source-hash: fnv1a64:d1b686b4a1391fd4
% canonical-lessons: JA-C470-kyoku, JA-C470-gakufu, JA-C470-min-you, JA-C470-enka, JA-C470-kuchibue

\chapter{Songs and Tunes}
\label{ch:ja-songs-and-tunes}
\hlchaptermodality{470}

\begin{chapteropening}
\textbf{By the end of this chapter:} \emph{I can talk about songs I like.}

Complete the last lesson of chapter 470: I can talk about songs I like.
\end{chapteropening}

\section[\texorpdfstring{\ja{きょく} (kyoku)}{きょく (kyoku)}]{\ja{きょく} (kyoku) — a piece of music; a tune}

\label{lesson:JA-C470-kyoku}



\begin{quote}
Before the new one: say the Japanese for to play music, then the Japanese for the tone of an instrument.
\end{quote}

\subsection*{You'll want to know: \ja{きょく}}
\textbf{\ja{きょく}} — \emph{kyoku} — \textquotedblleft{}a piece of music; a tune\textquotedblright{}.

\begin{quote}
\textbf{\ja{ははは} \ja{この} \ja{きょくが} \ja{すきです。}} — \emph{haha wa kono kyoku ga suki desu} — \textquotedblleft{}My mother likes this tune.\textquotedblright{}
\end{quote}

\begin{cousinweb}
\textbf{\ja{きょく}} (\emph{kyoku}) is Sino-Japanese (kango): one Chinese character read with its on'yomi, a reading borrowed from Chinese. It is \emph{kyoku}, \textquotedblleft{}melody\textquotedblright{} (Middle Chinese \emph{khjowk}) (KANJIDIC2; Middle Chinese in Baxter's transcription).
\end{cousinweb}

\begin{grammarlens}[title={What you've built}]
One more word for describing what you lived and felt.
\end{grammarlens}

\subsection*{Your turn}
\begin{itemize}
\raggedright
  \item \emph{Say it:} \emph{kyoku}
  \item \emph{Say it:} \emph{kyoku}, once more
  \item \emph{Say it:} \emph{haha wa kono kyoku ga suki desu} — My mother likes this tune.
  \item \emph{Recall:} say the Japanese for to play music, then the Japanese for the tone of an instrument, then say \emph{kyoku} again
\end{itemize}

\begin{tcolorbox}[breakable,colback=teal!4,colframe=teal!35!black,title={Before you move on}]
What does \ja{きょく} mean? (\textquotedblleft{}A piece of music; a tune\textquotedblright{}.) Say it once more.
\end{tcolorbox}
\section[\texorpdfstring{\ja{がくふ} (gakufu)}{がくふ (gakufu)}]{\ja{がくふ} (gakufu) — sheet music; a score}

\label{lesson:JA-C470-gakufu}



\begin{quote}
Before the new one: say the Japanese for the tone of an instrument, then the Japanese for a piece of music.
\end{quote}

\subsection*{You'll want to know: \ja{がくふ}}
\textbf{\ja{がくふ}} — \emph{gakufu} — \textquotedblleft{}sheet music; a score\textquotedblright{}.

\begin{quote}
\textbf{\ja{がくふを} \ja{よんで、うたいました。}} — \emph{gakufu o yonde, utaimashita} — \textquotedblleft{}I read the sheet music and sang.\textquotedblright{}
\end{quote}

\begin{cousinweb}
\textbf{\ja{がくふ}} (\emph{gakufu}) is Sino-Japanese (kango): two Chinese characters read with their on'yomi, readings borrowed from Chinese. They are \emph{gaku}, \textquotedblleft{}music\textquotedblright{} (Middle Chinese \emph{ngaewk}) and \emph{fu}, \textquotedblleft{}musical score\textquotedblright{} (\emph{puX}) (KANJIDIC2; Middle Chinese in Baxter's transcription).
\end{cousinweb}

\begin{grammarlens}[title={What you've built}]
One more word for describing what you lived and felt.
\end{grammarlens}

\subsection*{Your turn}
\begin{itemize}
\raggedright
  \item \emph{Say it:} \emph{gakufu}
  \item \emph{Say it:} \emph{gakufu}, once more
  \item \emph{Say it:} \emph{gakufu o yonde, utaimashita} — I read the sheet music and sang.
  \item \emph{Recall:} say the Japanese for the tone of an instrument, then the Japanese for a piece of music, then say \emph{gakufu} again
\end{itemize}

\begin{tcolorbox}[breakable,colback=teal!4,colframe=teal!35!black,title={Before you move on}]
What does \ja{がくふ} mean? (\textquotedblleft{}Sheet music; a score\textquotedblright{}.) Say it once more.
\end{tcolorbox}
\section[\texorpdfstring{\ja{みんよう} (min'yō)}{みんよう (min'yō)}]{\ja{みんよう} (min'yō) — a folk song}

\label{lesson:JA-C470-min-you}



\begin{quote}
Before the new one: say the Japanese for a piece of music, then the Japanese for sheet music.
\end{quote}

\subsection*{You'll want to know: \ja{みんよう}}
\textbf{\ja{みんよう}} — \emph{min'yō} — \textquotedblleft{}a folk song\textquotedblright{}.

\begin{quote}
\textbf{\ja{そぼは} \ja{みんようを} \ja{うたいます。}} — \emph{sobo wa min'yō o utaimasu} — \textquotedblleft{}My grandmother sings folk songs.\textquotedblright{}
\end{quote}

\begin{cousinweb}
\textbf{\ja{みんよう}} (\emph{min'yō}) is Sino-Japanese (kango): two Chinese characters read with their on'yomi, readings borrowed from Chinese. They are \emph{min}, \textquotedblleft{}people\textquotedblright{} (Middle Chinese \emph{mjin}) and \emph{yō}, \textquotedblleft{}song\textquotedblright{} (\emph{yew}) (KANJIDIC2; Middle Chinese in Baxter's transcription).
\end{cousinweb}

\begin{grammarlens}[title={What you've built}]
One more word for describing what you lived and felt.
\end{grammarlens}

\subsection*{Your turn}
\begin{itemize}
\raggedright
  \item \emph{Say it:} \emph{min'yō}
  \item \emph{Say it:} \emph{min'yō}, once more
  \item \emph{Say it:} \emph{sobo wa min'yō o utaimasu} — My grandmother sings folk songs.
  \item \emph{Recall:} say the Japanese for a piece of music, then the Japanese for sheet music, then say \emph{min'yō} again
\end{itemize}

\begin{tcolorbox}[breakable,colback=teal!4,colframe=teal!35!black,title={Before you move on}]
What does \ja{みんよう} mean? (\textquotedblleft{}A folk song\textquotedblright{}.) Say it once more.
\end{tcolorbox}
\section[\texorpdfstring{\ja{えんか} (enka)}{えんか (enka)}]{\ja{えんか} (enka) — an enka ballad}

\label{lesson:JA-C470-enka}



\begin{quote}
Before the new one: say the Japanese for sheet music, then the Japanese for a folk song.
\end{quote}

\subsection*{You'll want to know: \ja{えんか}}
\textbf{\ja{えんか}} — \emph{enka} — \textquotedblleft{}an enka ballad\textquotedblright{}.

\begin{quote}
\textbf{\ja{ちちは} \ja{えんかが} \ja{すきです。}} — \emph{chichi wa enka ga suki desu} — \textquotedblleft{}My father likes enka ballads.\textquotedblright{}
\end{quote}

\begin{cousinweb}
\textbf{\ja{えんか}} (\emph{enka}) is Sino-Japanese (kango): two Chinese characters read with their on'yomi, readings borrowed from Chinese. They are \emph{en}, \textquotedblleft{}performance\textquotedblright{} (Middle Chinese \emph{yenX}) and \emph{ka}, \textquotedblleft{}song\textquotedblright{} (\emph{ka}) (KANJIDIC2; Middle Chinese in Baxter's transcription).
\end{cousinweb}

\begin{grammarlens}[title={What you've built}]
One more word for describing what you lived and felt.
\end{grammarlens}

\subsection*{Your turn}
\begin{itemize}
\raggedright
  \item \emph{Say it:} \emph{enka}
  \item \emph{Say it:} \emph{enka}, once more
  \item \emph{Say it:} \emph{chichi wa enka ga suki desu} — My father likes enka ballads.
  \item \emph{Recall:} say the Japanese for sheet music, then the Japanese for a folk song, then say \emph{enka} again
\end{itemize}

\begin{tcolorbox}[breakable,colback=teal!4,colframe=teal!35!black,title={Before you move on}]
What does \ja{えんか} mean? (\textquotedblleft{}An enka ballad\textquotedblright{}.) Say it once more.
\end{tcolorbox}
\section[\texorpdfstring{\ja{くちぶえ} (kuchibue)}{くちぶえ (kuchibue)}]{\ja{くちぶえ} (kuchibue) — whistling; a whistle}

\label{lesson:JA-C470-kuchibue}



\begin{quote}
Before the new one: say the Japanese for a folk song, then the Japanese for an enka ballad.
\end{quote}

\subsection*{You'll want to know: \ja{くちぶえ}}
\textbf{\ja{くちぶえ}} — \emph{kuchibue} — \textquotedblleft{}whistling; a whistle\textquotedblright{}.

\begin{quote}
\textbf{\ja{かれは} \ja{くちぶえを} \ja{ふきました。}} — \emph{kare wa kuchibue o fukimashita} — \textquotedblleft{}He whistled.\textquotedblright{}
\end{quote}

\begin{cousinweb}
\textbf{\ja{くちぶえ}} (\emph{kuchibue}) is a native Japanese compound (wago), inherited rather than borrowed from Chinese: \emph{kuchi}, \textquotedblleft{}mouth\textquotedblright{} and \emph{bue}, \textquotedblleft{}flute (fue)\textquotedblright{} (JMdict, KANJIDIC2).
\end{cousinweb}

\begin{grammarlens}[title={What you've built}]
One more word for describing what you lived and felt.
\end{grammarlens}

\subsection*{Your turn}
\begin{itemize}
\raggedright
  \item \emph{Say it:} \emph{kuchibue}
  \item \emph{Say it:} \emph{kuchibue}, once more
  \item \emph{Say it:} \emph{kare wa kuchibue o fukimashita} — He whistled.
  \item \emph{Recall:} say the Japanese for a folk song, then the Japanese for an enka ballad, then say \emph{kuchibue} again
\end{itemize}

\begin{tcolorbox}[breakable,colback=teal!4,colframe=teal!35!black,title={Before you move on}]
What does \ja{くちぶえ} mean? (\textquotedblleft{}Whistling; a whistle\textquotedblright{}.) Say it once more.
\end{tcolorbox}
